\documentclass[UTF8,a4paper,11pt]{ctexart}

\usepackage{amsmath}
\usepackage{amssymb}
\usepackage{bm}
\usepackage{geometry}
\usepackage{hyperref}
\usepackage{listings}

\geometry{margin=2.5cm}
\hypersetup{colorlinks=true, linkcolor=blue, urlcolor=blue}

\lstset{
  basicstyle=\ttfamily\small,
  columns=fullflexible,
  breaklines=true,
  frame=single
}

\title{Structure-to-Surface Beam Mapper Comparison 中理论解的构造}
\author{}
\date{}

\begin{document}
\maketitle

\section{问题背景}

当前测试入口为
\begin{center}
\texttt{Bending\_Torsion\_Beam\_Test\_Case/run\_structure\_to\_surface\_mapper\_comparison.py}.
\end{center}
\noindent\emph{The current test entry point is
\texttt{Bending\_Torsion\_Beam\_Test\_Case/run\_structure\_to\_surface\_mapper\_comparison.py}.}

该脚本使用一个梁中心线模型作为 origin model part，使用包围梁的 surface mesh 作为 destination model part。
目标是比较以下三种映射结果在 surface mesh 上的准确度：
\par\noindent\emph{The script uses a beam centerline model as the origin model part and the surrounding surface mesh as the destination model part.
The purpose is to compare the accuracy of the following three mapping results on the surface mesh:}
\begin{itemize}
  \item \texttt{beam\_spline\_mapper};
  \item \texttt{beam\_mapper} with \texttt{use\_corotation = false};
  \item \texttt{beam\_mapper} with \texttt{use\_corotation = true}.
\end{itemize}

为了评价映射误差，需要先在每个 surface node 上构造一个解析参考位移，脚本中称为
\texttt{expected\_surface}。这个参考值不是由 mapper 得到，而是由梁中心线解析运动和截面有限转角直接计算得到。
\par\noindent\emph{To evaluate the mapping error, an analytical reference displacement must first be constructed at every surface node.
In the script this reference field is called \texttt{expected\_surface}.
It is not obtained from any mapper; instead, it is computed directly from the analytical beam-centerline motion and the finite rotation of the beam cross-section.}

\section{几何描述}

梁中心线位于初始构型中的全局 \(x\) 轴上：
\par\noindent\emph{In the initial configuration, the beam centerline lies on the global \(x\)-axis:}
\begin{equation}
  \bm{X}_c(x) =
  \begin{bmatrix}
    x \\ 0 \\ 0
  \end{bmatrix},
  \qquad x \in [0,L],
  \qquad L = 10.
\end{equation}

surface mesh 上任意一点的初始坐标记为
\par\noindent\emph{The initial coordinates of an arbitrary point on the surface mesh are denoted as}
\begin{equation}
  \bm{X}_s =
  \begin{bmatrix}
    x \\ y \\ z
  \end{bmatrix}.
\end{equation}

由于该 surface mesh 是围绕梁中心线生成的，点 \(\bm{X}_s\) 相对于同一 \(x\) 截面中心点的初始偏置向量为
\par\noindent\emph{Since the surface mesh is generated around the beam centerline, the initial offset vector of point \(\bm{X}_s\) with respect to the center point of the same \(x\)-section is}
\begin{equation}
  \bm{r}_0 =
  \bm{X}_s - \bm{X}_c(x)
  =
  \begin{bmatrix}
    0 \\ y \\ z
  \end{bmatrix}.
\end{equation}

\section{梁中心线解析运动}

\subsection{中心线弯曲角}

脚本沿用了原始 \texttt{beam\_geometry/ProjectParameters.json} 中给定的解析运动。定义
\par\noindent\emph{The script follows the analytical motion prescribed in the original \texttt{beam\_geometry/ProjectParameters.json}. Define}
\begin{equation}
  \theta(x) = 0.1 x \frac{\pi}{3}.
\end{equation}

令
\par\noindent\emph{Let}
\begin{equation}
  R_c = \frac{3L}{\pi} = \frac{30}{\pi}.
\end{equation}

这里的 \(R_c\) 可以理解为弯曲中心线的圆弧半径。由于
\par\noindent\emph{Here \(R_c\) can be interpreted as the radius of the circular arc of the bent centerline. Since}
\begin{equation}
  \frac{x}{R_c}
  =
  \frac{x}{30/\pi}
  =
  x\frac{\pi}{30}
  =
  0.1x\frac{\pi}{3}
  =
  \theta(x),
\end{equation}
所以 \(\theta(x)\) 正好是中心线圆弧参数角。
\par\noindent\emph{Therefore, \(\theta(x)\) is exactly the angular parameter of the centerline circular arc.}

\subsection{中心线位移}

原始直梁中心线坐标为 \((x,0,0)\)。弯曲后的中心线坐标取为
\par\noindent\emph{The original straight beam centerline coordinates are \((x,0,0)\). The coordinates of the bent centerline are taken as}
\begin{equation}
  \bm{x}_c(x) =
  \begin{bmatrix}
    R_c \sin \theta(x) \\
    -R_c + R_c \cos \theta(x) \\
    0
  \end{bmatrix}.
\end{equation}

因此中心线位移为
\par\noindent\emph{Therefore, the centerline displacement is}
\begin{equation}
  \bm{u}_c(x)
  =
  \bm{x}_c(x) - \bm{X}_c(x)
  =
  \begin{bmatrix}
    R_c \sin \theta(x) - x \\
    -R_c + R_c \cos \theta(x) \\
    0
  \end{bmatrix}.
\end{equation}

这对应脚本中的实现：
\par\noindent\emph{This corresponds to the following implementation in the script:}
\begin{lstlisting}[language=Python]
radius = 3.0 * 10.0 / math.pi
theta = 0.1 * x * (math.pi / 3.0)
displacement = (
    radius * math.sin(theta) - x,
    -radius + radius * math.cos(theta),
    0.0,
)
\end{lstlisting}

\subsection{梁截面转角向量}

在 Kratos 中，节点转角变量 \(\texttt{ROTATION}\) 以旋转向量形式给出：
\par\noindent\emph{In Kratos, the nodal rotation variable \(\texttt{ROTATION}\) is given in rotation-vector form:}
\begin{equation}
  \bm{\varphi}(x)
  =
  \begin{bmatrix}
    \theta(x) \\ 0 \\ -\theta(x)
  \end{bmatrix}.
\end{equation}

这对应脚本中的实现：
\par\noindent\emph{This corresponds to the following implementation in the script:}
\begin{lstlisting}[language=Python]
rotation = (theta, 0.0, -theta)
\end{lstlisting}

其中 \(x\) 分量表示绕全局 \(x\) 轴的扭转，\(z\) 分量表示与中心线弯曲相关的转动。
\par\noindent\emph{The \(x\)-component represents torsion around the global \(x\)-axis, while the \(z\)-component represents the rotation associated with the bending of the centerline.}

\section{由旋转向量构造有限转角旋转矩阵}

\subsection{旋转向量的轴角形式}

给定旋转向量
\par\noindent\emph{Given a rotation vector}
\begin{equation}
  \bm{\varphi} =
  \begin{bmatrix}
    \varphi_x \\ \varphi_y \\ \varphi_z
  \end{bmatrix},
\end{equation}
其模长为旋转角：
\par\noindent\emph{its magnitude is the rotation angle:}
\begin{equation}
  \alpha = \|\bm{\varphi}\|.
\end{equation}

若 \(\alpha \neq 0\)，单位旋转轴为
\par\noindent\emph{If \(\alpha \neq 0\), the unit rotation axis is}
\begin{equation}
  \bm{a} =
  \frac{\bm{\varphi}}{\alpha}
  =
  \begin{bmatrix}
    a_x \\ a_y \\ a_z
  \end{bmatrix}.
\end{equation}

\subsection{叉乘矩阵}

定义 \(\bm{a}\) 的反对称矩阵
\par\noindent\emph{Define the skew-symmetric matrix of \(\bm{a}\) as}
\begin{equation}
  [\bm{a}]_\times =
  \begin{bmatrix}
    0 & -a_z & a_y \\
    a_z & 0 & -a_x \\
    -a_y & a_x & 0
  \end{bmatrix}.
\end{equation}

对任意向量 \(\bm{v}\)，有
\par\noindent\emph{For any vector \(\bm{v}\),}
\begin{equation}
  [\bm{a}]_\times \bm{v} = \bm{a} \times \bm{v}.
\end{equation}

\subsection{Rodrigues 公式}

有限转角旋转矩阵由 Rodrigues 公式给出：
\par\noindent\emph{The finite-rotation matrix is given by Rodrigues' formula:}
\begin{equation}
  \bm{R}(\bm{\varphi})
  =
  \cos\alpha\,\bm{I}
  +
  (1-\cos\alpha)\,\bm{a}\otimes\bm{a}
  +
  \sin\alpha\,[\bm{a}]_\times.
\end{equation}

当 \(\alpha\) 足够小时，取
\par\noindent\emph{When \(\alpha\) is sufficiently small, we use}
\begin{equation}
  \bm{R}(\bm{\varphi}) = \bm{I}.
\end{equation}

这对应脚本中的实现：
\par\noindent\emph{This corresponds to the following implementation in the script:}
\begin{lstlisting}[language=Python]
def rotation_matrix_from_rotation_vector(rotation):
    angle = math.sqrt(sum(value * value for value in rotation))
    if angle < 1e-14:
        return (
            (1.0, 0.0, 0.0),
            (0.0, 1.0, 0.0),
            (0.0, 0.0, 1.0),
        )

    axis = tuple(value / angle for value in rotation)
    axis_cross = (
        (0.0, -axis[2], axis[1]),
        (axis[2], 0.0, -axis[0]),
        (-axis[1], axis[0], 0.0),
    )
    axis_outer = tuple(tuple(axis[i] * axis[j] for j in range(3)) for i in range(3))
    identity = (
        (1.0, 0.0, 0.0),
        (0.0, 1.0, 0.0),
        (0.0, 0.0, 1.0),
    )
    cos_angle = math.cos(angle)
    sin_angle = math.sin(angle)

    return tuple(
        tuple(
            cos_angle * identity[i][j]
            + (1.0 - cos_angle) * axis_outer[i][j]
            + sin_angle * axis_cross[i][j]
            for j in range(3)
        )
        for i in range(3)
    )
\end{lstlisting}

\section{Surface 点理论位移}

\subsection{有限转角截面运动}

surface 点的初始位置可以写成
\par\noindent\emph{The initial position of a surface point can be written as}
\begin{equation}
  \bm{X}_s = \bm{X}_c(x) + \bm{r}_0.
\end{equation}

假设该 surface 点随梁截面作刚体运动，则变形后位置为
\par\noindent\emph{Assuming that the surface point follows the beam cross-section as a rigid body, its deformed position is}
\begin{equation}
  \bm{x}_s
  =
  \bm{x}_c(x)
  +
  \bm{R}\bigl(\bm{\varphi}(x)\bigr)\bm{r}_0.
\end{equation}

因此 surface 点理论位移为
\par\noindent\emph{Therefore, the theoretical displacement of the surface point is}
\begin{align}
  \bm{u}_s^{\mathrm{ref}}
  &=
  \bm{x}_s - \bm{X}_s \\
  &=
  \left(\bm{x}_c(x)-\bm{X}_c(x)\right)
  +
  \left(\bm{R}\bigl(\bm{\varphi}(x)\bigr)\bm{r}_0-\bm{r}_0\right) \\
  &=
  \bm{u}_c(x)
  +
  \left(\bm{R}\bigl(\bm{\varphi}(x)\bigr)-\bm{I}\right)\bm{r}_0.
\end{align}

这就是脚本中 \texttt{expected\_surface} 的理论解。
\par\noindent\emph{This is the theoretical solution used for \texttt{expected\_surface} in the script.}

\subsection{与小转角公式的关系}

若转角很小，则
\par\noindent\emph{If the rotation is small, then}
\begin{equation}
  \bm{R}(\bm{\varphi})
  \approx
  \bm{I} + [\bm{\varphi}]_\times.
\end{equation}

此时
\par\noindent\emph{In this case,}
\begin{equation}
  \bm{u}_s^{\mathrm{ref}}
  \approx
  \bm{u}_c(x)
  +
  \bm{\varphi}(x) \times \bm{r}_0.
\end{equation}

这就是 linear 版本 mapper 常见的小转角近似。但当前为了评估 \texttt{use\_corotation=true} 的 \texttt{beam\_mapper}，脚本采用的是有限转角表达式
\par\noindent\emph{This is the small-rotation approximation commonly used by the linear version of the mapper.
However, in order to evaluate \texttt{beam\_mapper} with \texttt{use\_corotation=true}, the current script uses the finite-rotation expression}
\begin{equation}
  \bm{u}_c(x)
  +
  \left(\bm{R}(\bm{\varphi}(x))-\bm{I}\right)\bm{r}_0,
\end{equation}
而不是小转角表达式。
\par\noindent\emph{rather than the small-rotation expression.}

\section{编程实现}

对每个 surface node，脚本执行以下步骤：
\par\noindent\emph{For each surface node, the script performs the following steps:}
\begin{enumerate}
  \item 读取节点初始坐标 \((x,y,z)\)；
  \par\noindent\emph{Read the initial nodal coordinates \((x,y,z)\).}
  \item 用 \(x\) 计算中心线位移 \(\bm{u}_c(x)\) 和旋转向量 \(\bm{\varphi}(x)\)；
  \par\noindent\emph{Use \(x\) to compute the centerline displacement \(\bm{u}_c(x)\) and the rotation vector \(\bm{\varphi}(x)\).}
  \item 构造 offset 向量
  \[
    \bm{r}_0 = (0,y,z)^T;
  \]
  \par\noindent\emph{Construct the offset vector \(\bm{r}_0 = (0,y,z)^T\).}
  \item 由 Rodrigues 公式构造旋转矩阵 \(\bm{R}(\bm{\varphi})\)；
  \par\noindent\emph{Construct the rotation matrix \(\bm{R}(\bm{\varphi})\) using Rodrigues' formula.}
  \item 计算旋转后的 offset：
  \[
    \bm{r} = \bm{R}(\bm{\varphi})\bm{r}_0;
  \]
  \par\noindent\emph{Compute the rotated offset \(\bm{r} = \bm{R}(\bm{\varphi})\bm{r}_0\).}
  \item 得到 surface 理论位移：
  \[
    \bm{u}_s^{\mathrm{ref}} = \bm{u}_c(x) + \bm{r} - \bm{r}_0;
  \]
  \par\noindent\emph{Obtain the theoretical surface displacement
  \(\bm{u}_s^{\mathrm{ref}} = \bm{u}_c(x) + \bm{r} - \bm{r}_0\).}
  \item 将该值写入 \texttt{expected\_surface} 的 \texttt{DISPLACEMENT}。
  \par\noindent\emph{Write this value into the \texttt{DISPLACEMENT} field of \texttt{expected\_surface}.}
\end{enumerate}

对应代码为：
\par\noindent\emph{The corresponding code is:}
\begin{lstlisting}[language=Python]
def fill_expected_surface_displacement(surface):
    for node in surface.Nodes:
        center_displacement, rotation = prescribed_beam_kinematics(node.X0)
        offset = (0.0, node.Y0, node.Z0)
        rotation_matrix = rotation_matrix_from_rotation_vector(rotation)
        rotated_offset = mat_vec_prod(rotation_matrix, offset)
        rotational_displacement = tuple(rotated_offset[i] - offset[i] for i in range(3))
        node.SetSolutionStepValue(
            KM.DISPLACEMENT,
            [center_displacement[i] + rotational_displacement[i] for i in range(3)],
        )
\end{lstlisting}

\section{误差度量}

对任意 mapper 得到的 surface 位移 \(\bm{u}_s^{\mathrm{map}}\)，误差定义为欧氏范数：
\par\noindent\emph{For the surface displacement \(\bm{u}_s^{\mathrm{map}}\) obtained by any mapper, the error is defined as the Euclidean norm:}
\begin{equation}
  e =
  \left\|
    \bm{u}_s^{\mathrm{map}}
    -
    \bm{u}_s^{\mathrm{ref}}
  \right\|_2.
\end{equation}

在代码中：
\par\noindent\emph{In the code:}
\begin{lstlisting}[language=Python]
def norm_3d(lhs, rhs):
    return math.sqrt(sum((lhs[i] - rhs[i]) ** 2 for i in range(3)))
\end{lstlisting}

三个 mapper 的最大误差分别为
\par\noindent\emph{The maximum errors of the three mappers are respectively}
\begin{align}
  e_{\max}^{\mathrm{spline}}
  &=
  \max_i
  \left\|
    \bm{u}_{s,i}^{\mathrm{beam\_spline}}
    -
    \bm{u}_{s,i}^{\mathrm{ref}}
  \right\|_2, \\
  e_{\max}^{\mathrm{linear}}
  &=
  \max_i
  \left\|
    \bm{u}_{s,i}^{\mathrm{beam\_mapper,linear}}
    -
    \bm{u}_{s,i}^{\mathrm{ref}}
  \right\|_2, \\
  e_{\max}^{\mathrm{corotation}}
  &=
  \max_i
  \left\|
    \bm{u}_{s,i}^{\mathrm{beam\_mapper,corotation}}
    -
    \bm{u}_{s,i}^{\mathrm{ref}}
  \right\|_2.
\end{align}

\section{为什么 co-rotation 版本应更接近理论解}

当前理论解使用的是有限转角截面刚体运动：
\par\noindent\emph{The current theoretical solution uses finite-rotation rigid-body motion of the cross-section:}
\begin{equation}
  \bm{u}_s^{\mathrm{ref}}
  =
  \bm{u}_c(x)
  +
  \left(\bm{R}(\bm{\varphi})-\bm{I}\right)\bm{r}_0.
\end{equation}

\texttt{beam\_mapper} 的 \texttt{use\_corotation=true} 分支同样试图在 co-rotational 框架中处理截面有限转动，因此它应当比 \texttt{use\_corotation=false} 的小转角线性版本更接近该理论解。
\par\noindent\emph{The \texttt{use\_corotation=true} branch of \texttt{beam\_mapper} also attempts to handle finite cross-section rotations in a co-rotational frame.
Therefore, it should be closer to this theoretical solution than the small-rotation linear version with \texttt{use\_corotation=false}.}

如果使用小转角参考
\par\noindent\emph{If the small-rotation reference is used}
\begin{equation}
  \bm{u}_s^{\mathrm{small}}
  =
  \bm{u}_c(x)
  +
  \bm{\varphi}\times\bm{r}_0,
\end{equation}
则 linear mapper 的误差会相对更小；但这并不能检验 co-rotation 对有限转角的优势。因此当前脚本选择有限转角理论解作为主要参考。
\par\noindent\emph{then the error of the linear mapper would become relatively smaller.
However, that would not test the advantage of co-rotation for finite rotations.
For this reason, the current script uses the finite-rotation theoretical solution as the main reference.}

\section{当前脚本输出}

脚本会输出如下 VTK 文件：
\par\noindent\emph{The script writes the following VTK files:}
\begin{itemize}
  \item \texttt{origin\_beam\_structure\_Parts\_Beam\_beam\_0\_0.vtk};
  \item \texttt{expected\_surface\_Parts\_Shell\_wet\_surface\_0\_0.vtk};
  \item \texttt{beam\_spline\_mapped\_surface\_Parts\_Shell\_wet\_surface\_0\_0.vtk};
  \item \texttt{beam\_mapper\_linear\_mapped\_surface\_Parts\_Shell\_wet\_surface\_0\_0.vtk};
  \item \texttt{beam\_mapper\_corotation\_mapped\_surface\_Parts\_Shell\_wet\_surface\_0\_0.vtk}.
\end{itemize}

输出目录为
\begin{center}
\texttt{Bending\_Torsion\_Beam\_Test\_Case/TestCase\_Output/structure\_to\_surface\_10x100}.
\end{center}
\noindent\emph{The output directory is
\texttt{Bending\_Torsion\_Beam\_Test\_Case/TestCase\_Output/structure\_to\_surface\_10x100}.}

\end{document}
